%% kapa: exam=1 round=29 year=2018 period=2 form=A subject=accounting subject_ko=회계학 seq=8 number=48
%% source: Q-Net 감정평가사 기출문제 (cst00309) artlSeq=5202808 file=extracted/1008/5202808/2018년 제29회 감정평가사 1차시험 기출문제/2018년 제29회 감정평가사 1차시험 기출문제/2교시 A형.pdf
%% conversion: kordoc 4.24.51 -> tools/gen1.py (auto); review=proofread-sonnet; flags=none
\begin{question}{48}
(주)감평은 20x8년 3월 1일 사용중이던 기계장치를 (주)대한의 신형 기계장치와 교환하면서 ￦4,000의 현금을 추가로 지급하였다. (주)감평이 사용하던 기계장치는 20x5년에 ￦41,000에 취득한 것으로 교환당시 감가상각누계액은 ￦23,000이고 공정가치는 ￦21,000이다. 한편, 교환시점 (주)대한의 신형 기계장치의 공정가치는 ￦26,000이다. 동 교환거래가 상업적 실질이 있으며 (주)감평의 사용중이던 기계장치의 공정가치가 더 명백한 경우 (주)감평이 교환거래로 인해 인식할 처분손익은?
\begin{choices}
  \item 이익 ￦3,000
  \item 이익 ￦4,000
  \item 손실 ￦3,000
  \item 손실 ￦4,000
  \item 이익 ￦1,000
\end{choices}
\end{question}
\answer{1}
